\section{Exact comparisons with one sign test}
\label{sec:upper}

Branch-and-bound methods for mixed-integer nonlinear programs make exact
decisions. A node is pruned when its relaxation value is at least the
incumbent value. Bound activity at a relaxation optimizer informs branching
decisions. Two subproblems are ranked by their optimal values. Approximate
values and points do not settle such decisions by themselves. For the
classes studied here, a nonzero difference can be as small as
$2^{-2^{\poly(L)}}$, so a sign certified by approximation can require
exponentially many printed bits.

For globally strongly convex explicit polynomials, this section reduces every
order or equality comparison of a polynomial observable at the minimizer to
the sign of one integer arithmetic circuit, that is, to one instance of
$\PosSLP$. The circuit is built in polynomial time without oracle calls. The lower bounds of
Section~\ref{sec:reductions} show that the order comparisons are also
$\PosSLP$-hard for certified quartics. In that class the order comparisons
are therefore equivalent to $\PosSLP$ under polynomial-time many-one
reductions.

The proof has four parts. An ordinary algorithm computes a starting point
with polynomially many bits. Newton's method, run as a shared rational
circuit, reaches doubly exponential accuracy with polynomially many gates. A
real-algebraic separation bound, used only in the analysis, determines how
much accuracy is needed. Finally, the observable is shifted by half of the
separation gap and division is removed without changing signs. The same
method gives an upper bound for zeros of strongly monotone maps that are not
gradients, for polyhedra on which a gap between zero and nonzero optimal
slacks is supplied, and, with an elementary separation bound, for the
certified cube-root circuits of Section~\ref{sec:reductions}. The last
subsection compiles an adaptive polynomial-time computation with exact sign
tests into a single sign test.

\subsection{Conventions}
\label{sec:upper-conventions}

We use the input and circuit models of Section~\ref{sec:models}. The
following points are used repeatedly.

\begin{enumerate}[label=(\roman*)]
\item \emph{Input.} Rational numbers are written in binary. An explicit
sparse polynomial is the list of its nonzero coefficients together with their
exponent vectors, written in unary, and the dimension $n$ is written in unary
as well; a polynomial map
$T=(T_1,\ldots,T_n)$ is the list of its components. We write $L\ge2$ for the
total binary length of the input, including the curvature bound, the
observable, thresholds, constraints and certificates. Thus every printed
numerator and denominator has absolute value at most $2^L$, every listed
polynomial has at most $L$ terms, and the dimension $n$ and the numerical
degree $D$ are at most $L$. If an encoding only guarantees
$D\le\poly(L)$, the results hold after replacing $L$ by that polynomial.
Sparse polynomials with binary exponents and no polynomial degree bound are
not covered: evaluating them at a short rational point can already produce
exponentially many bits.
\item \emph{Circuits.} An \emph{integer circuit} is a directed acyclic graph
whose sources are the constants $0$ and $1$ and whose other nodes are binary
$+$, $-$ and $\times$ gates; its size is its number of nodes. A
\emph{rational circuit} also allows rational constants written in binary and
division by nonzero values. A circuit for a vector has several output nodes,
and all nodes may be shared. Circuit values can have exponentially many bits;
they are never expanded. $\PosSLP$ asks whether the output of an integer
circuit is positive.
\item \emph{Reductions.} A problem \emph{reduces to one $\PosSLP$
instance} if a deterministic polynomial-time algorithm, making no oracle
calls, maps every input to an integer circuit whose output is positive
exactly when the answer is yes. For a promise problem this is required only on
inputs satisfying the promise. A supplied curvature bound is a promise and is
not checked. The certificate formats below are checked in polynomial time;
an invalid certificate is mapped to the circuit $0$, so these formats define
ordinary languages.
\item \emph{Comparisons.} We write ${\bowtie}$ for one of the relations
$>,\ge,<,\le,=,\ne$. A rational threshold is absorbed into the observable:
the value comparison $\min f\bowtie r$ uses $h=f-r$, and the coordinate
comparison $p_j\bowtie r$ uses $h=X_j-r$.
\end{enumerate}

\subsection{Three tools}
\label{sec:upper-tools}

The first tool says how close to zero a nonzero algebraic quantity defined
by a short real formula can be. It is used only in proofs; the reductions
never perform quantifier elimination. For an integer $N$, its bit size is
$\lceil\log_2(|N|+1)\rceil$.

\begin{lemma}[Algebraic separation]
\label{lem:separation}
There is an absolute integer constant $c_0\ge1$ with the following property.
Let $\Phi(y,z)$ be a quantifier-free formula in real variables $y\in\R^s$,
$s\ge1$, and $z\in\R$, that is, a Boolean combination of conditions
$P(y,z)>0$, $P(y,z)=0$ and $P(y,z)<0$, where every polynomial
$P\in\Z[y,z]$ that occurs has total degree at most $d\ge2$ and coefficients
of bit size at most $\tau\ge1$. Suppose that the set
$\{z\in\R:\exists y\in\R^s\ \Phi(y,z)\}$ consists of a single point
$\alpha$. Then $\alpha$ is algebraic, and the following hold.
\begin{enumerate}[label=(\alph*)]
\item \emph{Polynomial-size formulas.} If $s,d,\tau\le L^c$ for integers
$L\ge2$ and $c\ge1$, then $\alpha=0$ or $|\alpha|\ge2^{-2^{e}}$ with
$e=c(c_0+2)L^{c+1}$.
\item \emph{Explicit form.} If $\alpha\ne0$, then
$|\alpha|\ge2^{-2\tau d^{c_0s}}$. For fixed $d$, the exponent is linear in
$\tau$, with a factor that depends only on $s$.
\end{enumerate}
\end{lemma}

The proof in Appendix~\ref{app:upper-separation} applies a one-block
quantifier-elimination theorem with bounds on degrees and coefficient bit
sizes \citep[Theorem~2.27]{Basu2014}. Some polynomial of the eliminated
formula must vanish at $\alpha$, because otherwise the formula would define a
neighborhood of $\alpha$; an elementary root bound then gives the estimate.
The lemma does not require the defining system to have finitely many complex
solutions. For instance, $f(x,y)=(x^2+y^2)^2+x^2+y^2$ satisfies
$\nabla^2f\succeq2I$ on $\R^2$, yet its complex critical locus contains the
curve $x^2+y^2=-1/2$. Only the real projection matters. Form~(a) is used
for problems with polynomially many variables. Form~(b) keeps the dependence
on $\tau$ linear and the dependence on $s$ single exponential for fixed $d$,
which is what the parameterized algorithms of Section~\ref{sec:constraints}
need.

The second tool removes division from a rational circuit. The rule
$(N_a/Q_a)/(N_b/Q_b)=(N_aQ_b)/(Q_aN_b)$ can produce a negative denominator,
which would reverse the sign of the final numerator; detecting this needs a
sign test. Squaring the divisor's numerator avoids the problem.

\begin{lemma}[Division elimination with positive denominators]
\label{lem:denominator-clearing}
Let $C$ be a rational circuit with $\sigma$ nodes whose rational constants
have total bit length $b$. In time polynomial in $\sigma+b$, and without
sign tests, one can construct an integer circuit of size $O(\sigma+b)$ that
contains, for every node $v$ of $C$, two nodes $N_v,Q_v$ with the following
property. If every division gate of $C$ divides by a nonzero value, then
$Q_v>0$ and the value of $v$ equals $N_v/Q_v$ for every node $v$. In that
case the value of $v$ is positive if and only if $N_v>0$, so the sign of any
node of $C$ is decided by one $\PosSLP$ instance.
\end{lemma}

Tracking numerators and denominators is standard in $\PosSLP$ reductions
\citep{AllenderEtAl2009,EtessamiStewartYannakakis2012}. Division by a node
with numerator $N_b$ and denominator $Q_b$ is implemented as
\begin{equation}
 \frac{N_a/Q_a}{N_b/Q_b}=\frac{N_aQ_bN_b}{Q_aN_b^2},
 \label{eq:upper-division}
\end{equation}
whose denominator is positive whenever $Q_a>0$ and $N_b\ne0$. The full
proof is in Appendix~\ref{app:upper-clearing}. We state the lemma once
because every result in this section and several in
Sections~\ref{sec:constraints} and~\ref{sec:heights} use it.

The third tool is a short circuit for a very accurate approximation of the
optimizer. It is uniform in the observable: one circuit serves every
observable up to a stated encoding length. This uniformity is used when many
observables are tested at the same point, as in the active-set certificates
of Section~\ref{sec:constraints} and the circuit Grams of
Theorem~\ref{thm:circuit-gram}.

\begin{lemma}[Shared Newton point]
\label{lem:newton-circuit}
Let an input of length $L$ consist of $n\ge1$, a rational $\mu>0$, and
either
\begin{enumerate}[label=(\alph*)]
\item an explicit sparse polynomial $f\in\Q[X_1,\ldots,X_n]$ with the
promise $\nabla^2f(x)\succeq\mu I$ for all $x\in\R^n$, in which case put
$T=\nabla f$; or
\item an explicit sparse polynomial map $T\colon\R^n\to\R^n$ with the promise
$\ip{T(x)-T(y)}{x-y}\ge\mu\norm{x-y}^2$ for all $x,y\in\R^n$.
\end{enumerate}
Then $T$ has a unique zero $p$. Given in addition an integer $L'\ge L$,
written in unary, one can construct in deterministic time polynomial in
$L'$, without oracle calls, a shared rational circuit with outputs
$\hat x\in\Q^n$ in which every division is by a nonzero value, such that the
following holds for every explicit sparse polynomial
$h\in\Q[X_1,\ldots,X_n]$ of encoding length at most $L'$:
\[
 h(p)=0\ \text{ or }\ |h(p)|\ge g_{L'},
 \qquad
 |h(\hat x)-h(p)|\le g_{L'}/8,
 \qquad
 g_{L'}:=2^{-2^{e(L')}},
\]
where $e(L'):=(c_0+3)L'^2$ and $c_0$ is the constant of
Lemma~\ref{lem:separation}.
\end{lemma}

The circuit consists of an ordinary rational starting point $x_0$ with
$\norm{x_0-p}\le\mu/(4\Lambda)$, where $\Lambda=2^{10L'^2}$ bounds the
relevant derivatives, followed by $e(L')+1$ exact Newton steps
\[
 x_{t+1}=x_t-J_T(x_t)^{-1}T(x_t).
\]
In case~(a), the starting point comes from polynomial-time convex value
optimization (Lemma~\ref{lem:convex-value}) and strong convexity. In
case~(b) no objective is available, and we use the ellipsoid method with the
cuts $\{y:T(x)^{\mathsf T}(y-x)\le0\}$; every such cut keeps a fixed ball
around $p$, which forces the method to reach a point with small residual
(Lemma~\ref{lem:upper-warm-start}). The Newton error contracts as
$\norm{x_{t+1}-p}\le(\Lambda/2\mu)\norm{x_t-p}^2$, so after $e(L')+1$ steps
it is at most $2^{-6\cdot2^{e(L')}}$. Each step solves a linear system whose
matrix has a positive definite symmetric part; Gaussian elimination without
pivoting has positive pivots, so no sign test is needed. Sharing the nodes of
earlier steps keeps the size polynomial, although the expanded coordinates
of $\hat x$ can have exponentially many bits. The proof is in
Appendices~\ref{app:upper-newton} and~\ref{app:upper-warm}.

\subsection{Strongly convex objectives}
\label{sec:upper-convex}

\begin{theorem}[Exact comparison at a strongly convex minimizer]
\label{thm:exact-upper}
Let the input consist of $n\ge1$, explicit sparse polynomials
$f,h\in\Q[X_1,\ldots,X_n]$, a relation ${\bowtie}$, and one of the following.
\begin{enumerate}[label=(\alph*)]
\item A rational $\mu>0$, with the promise $\nabla^2f(x)\succeq\mu I$ for
every $x\in\R^n$.
\item A Hessian certificate: an explicit list $w(X,v)=(w_1,\ldots,w_m)$ of
distinct monomials of the form $X^\beta v_i$ that contains each of
$v_1,\ldots,v_n$, and a rational symmetric $m\times m$ matrix $M$. The
certificate is valid if $M\succ0$ and
$v^{\mathsf T}\nabla^2f(X)\,v=w(X,v)^{\mathsf T}Mw(X,v)$ as polynomials in
$(X,v)$.
\end{enumerate}
Let $p$ be the unique minimizer of $f$ on $\R^n$. Deciding
$h(p)\bowtie0$ reduces to one $\PosSLP$ instance: in case~(a) on all
promised inputs, and in case~(b) on all inputs, where the language consists of
the inputs with a valid certificate and $h(p)\bowtie0$. No lower bound on
$|h(p)|$ is assumed.
\end{theorem}

\begin{proof}[Proof outline]
Apply Lemma~\ref{lem:newton-circuit}(a) with $L'=L$ and put
$\hat\alpha=h(\hat x)$ and $g=g_L$, computed by $e(L)$ squarings of $1/2$.
Since $\alpha=h(p)$ is zero or at least $g$ in absolute value, and
$|\hat\alpha-\alpha|\le g/8$, each rational circuit in the following table
is positive exactly when its relation holds. Its absolute value is at least
$3g/8$ in the four order rows and at least $15g^2/64$ in the two equality
rows, so it is never zero.
\begin{center}
\begin{tabular}{ll@{\qquad}ll}
\toprule
Relation & Circuit & Relation & Circuit\\
\midrule
$\alpha>0$ & $\hat\alpha-g/2$ & $\alpha<0$ & $-\hat\alpha-g/2$\\
$\alpha\ge0$ & $\hat\alpha+g/2$ & $\alpha\le0$ & $-\hat\alpha+g/2$\\
$\alpha=0$ & $g^2/4-\hat\alpha^2$ & $\alpha\ne0$ & $\hat\alpha^2-g^2/4$\\
\bottomrule
\end{tabular}
\end{center}
Lemma~\ref{lem:denominator-clearing} turns the selected rational circuit into
an integer circuit with the same sign. For case~(b), the identity is checked
by sparse coefficient comparison, $M\succ0$ by rational symmetric
elimination, and then $\mu=\det M/(\tr M)^{m-1}$ is a rational lower bound
for the least eigenvalue of $M$ (Lemma~\ref{lem:models-det-trace}). Since $w$
contains every $v_i$, $v^{\mathsf T}\nabla^2f(X)v\ge\mu\norm{w}^2\ge\mu\norm v^2$.
Case~(a) then applies with $L$ replaced by the polynomially larger length of
the input together with $\mu$. Details are in
Appendix~\ref{app:upper-convex-proof}.
\end{proof}

The theorem decides the sign of $\min f-r$ and of $p_j-r$, including the
equality cases, with no promise that these numbers are nonzero. For quartics
with the full basis $w=(v,X\otimes v)$, the strict and weak comparisons are
$\PosSLP$-hard (Section~\ref{sec:reductions}), and
Theorem~\ref{thm:quartic-complete} records the resulting classification.
Equality has only the upper bound: none of the lower reductions in this paper
produces equality instances. The certificate in~(b) is a sufficient format,
not a recognition algorithm for strong convexity. Degree four is the first
degree at which the theorem has content: a globally convex polynomial of
degree at most three is quadratic (Lemma~\ref{lem:models-low-degree}), and
then $p$ and $h(p)$ are obtained by rational linear algebra.

For certified quartics, nonnegativity of $f$ and membership in the cone of
real sums of squares both mean $\min f\ge0$, and the existence of a positive
definite polynomial Gram matrix means $\min f>0$; these equivalences are
recorded with Theorem~\ref{thm:quartic-complete}, and Theorem~\ref{thm:circuit-gram}
gives the corresponding circuit Gram. The three certificate questions
therefore inherit the one-instance upper bound. Rational sum-of-squares
membership is different: when $\min f=0$ it is not determined by the sign of
the minimum (Section~\ref{sec:fields}), and no matching upper bound is claimed
for it.

\begin{corollary}[Comparing two minima]
\label{cor:upper-two-minima}
Let $f_1\in\Q[X_1,\ldots,X_{n_1}]$ and $f_2\in\Q[Y_1,\ldots,Y_{n_2}]$ each
satisfy the hypotheses of Theorem~\ref{thm:exact-upper}, both with
promises~(a) or both with certificates~(b). For each relation ${\bowtie}$,
deciding $\min f_1\bowtie\min f_2$ reduces to one $\PosSLP$ instance.
\end{corollary}

\begin{proof}
Use disjoint variables and put $F(X,Y)=f_1(X)+f_2(Y)$ and
$h(X,Y)=f_1(X)-f_2(Y)$. The Hessian of $F$ is block diagonal, so
$\nabla^2F\succeq\min\{\mu_1,\mu_2\}I$. With certificates, check both, map
invalid inputs to the circuit $0$, and take
$\mu_i=\det M_i/(\tr M_i)^{m_i-1}$. The minimizer of $F$ is the pair of minimizers,
and there $h$ equals $\min f_1-\min f_2$. Apply
Theorem~\ref{thm:exact-upper}(a) to $F$ and $h$.
\end{proof}

The proof needs only a curvature bound for the sum. A block-separated sum
need not have a positive definite Hessian Gram on the full joint basis, and
none is required.

\begin{proposition}[Circuit witnesses for strict inequalities]
\label{prop:upper-circuit-witness}
Under the hypotheses of Theorem~\ref{thm:exact-upper}(a), the point $\hat x$
of Lemma~\ref{lem:newton-circuit} with $L'=L$ is a rational vector with a
shared circuit of size polynomial in $L$, constructed without oracle calls
and without knowing any sign. If $h(p)>0$, then $h(\hat x)>0$. In particular,
if $\min f<0$ then $f(\hat x)<0$.
\end{proposition}

\begin{proof}
If $h(p)>0$, then $h(p)\ge g_L$ and $h(\hat x)\ge7g_L/8>0$. Apply this to
$h=-f$ for the last claim.
\end{proof}

Thus every strictly feasible instance $\{f<0\}\ne\emptyset$ has a short
feasible witness in the shared-circuit format, built by the same algorithm
on every input. Three limits matter. Checking that $f(\hat x)<0$ is again a
$\PosSLP$ question. The expanded numerators and denominators of $\hat x$ can
be exponentially long, and Section~\ref{sec:heights} shows that for some
strictly feasible quartics every rational feasible point has exponentially
many bits. When $\min f=0$, the set $\{f\le0\}$ is the single point $p$,
which can be irrational (Theorem~\ref{thm:singleton-field}); then no rational
witness exists at all.

\subsection{Zeros of strongly monotone polynomial maps}
\label{sec:upper-monotone}

The same comparison holds at the zero of a polynomial map that is not a
gradient. Such maps describe equilibria of games and variational models
without a potential. The only new ingredient is the starting point, because
there is no objective whose value can be minimized.

\begin{theorem}[Zeros of strongly monotone maps]
\label{thm:upper-monotone}
Let the input consist of $n\ge1$, an explicit sparse polynomial map
$T=(T_1,\ldots,T_n)$, an explicit sparse polynomial $h$, a relation
${\bowtie}$, and one of the following.
\begin{enumerate}[label=(\alph*)]
\item A rational $\mu>0$, with the promise
$\ip{T(x)-T(y)}{x-y}\ge\mu\norm{x-y}^2$ for all $x,y\in\R^n$.
\item A certificate $(w,M)$ as in Theorem~\ref{thm:exact-upper}(b), valid if
$M\succ0$ and $v^{\mathsf T}J_T(X)\,v=w(X,v)^{\mathsf T}Mw(X,v)$ as
polynomials.
\end{enumerate}
Then $T$ has a unique zero $p$, and deciding $h(p)\bowtie0$ reduces to one
$\PosSLP$ instance, in the same sense as in Theorem~\ref{thm:exact-upper}.
\end{theorem}

The Jacobian $J_T$ need not be symmetric; the certificate controls only its
symmetric part, which suffices because
$\ip{T(x)-T(y)}{x-y}=\int_0^1(x-y)^{\mathsf T}J_T(y+\theta(x-y))(x-y)\,d\theta$.
Existence of the zero is proved, not assumed. The proof
(Appendix~\ref{app:upper-monotone-proof}) uses Lemma~\ref{lem:newton-circuit}(b);
the Newton steps solve the nonsymmetric system directly by elimination
without pivoting.

\begin{example}[Maps that are not gradients]
\label{ex:upper-nongradient}
For a rational skew-symmetric matrix $A$ and rational $b$, the cubic map
$T(x)=(1+\norm x^2)x+Ax-b$ has
$J_T(x)=(1+\norm x^2)I+2xx^{\mathsf T}+A$, hence
$v^{\mathsf T}J_T(x)v=\norm v^2+\norm{x\otimes v}^2+2(x^{\mathsf T}v)^2$.
A positive definite certificate on $(v,x\otimes v)$ is $I\oplus(I+2\,dd^{\mathsf T})$,
where $d$ selects the coordinates $x_iv_i$. If $A\ne0$, then $J_T$ is not
symmetric and $T$ is not a gradient. The antisymmetric part of the Jacobian
can also vary with $x$. In two variables, let
\[
 T(x,y)=(1+x^2+y^2)\,(x,y)+(x^2y/2,\,0).
\]
With $z=(xv_1,xv_2,yv_1,yv_2)$ one checks
$v^{\mathsf T}J_T v=\norm v^2+z^{\mathsf T}Q_0z$ for
\[
 Q_0=\begin{pmatrix}3&1/4&1/2&2\\1/4&1&0&0\\1/2&0&1&0\\2&0&0&3\end{pmatrix},
\]
whose leading principal minors are $3$, $47/16$, $43/16$ and $65/16$. Here
$(J_T)_{12}-(J_T)_{21}=x^2/2$, so $T$ is not the sum of a gradient and a
constant linear map.
\end{example}

For cubic maps with the full basis $w=(v,X\otimes v)$, the gradients of the
certified quartics of Section~\ref{sec:reductions} are instances of
Theorem~\ref{thm:upper-monotone}(b). Hence strict and weak comparisons at
zeros of certified cubic maps are $\PosSLP$-complete, while equality again
has only the upper bound. Cubic degree is the first degree with content: a
strongly monotone polynomial map of degree at most two is affine
(Lemma~\ref{lem:models-low-degree}), and its zero is rational. Constrained
variational inequalities with such maps are treated in
Section~\ref{sec:constraints}.

\subsection{Polyhedra with a supplied slack gap}
\label{sec:upper-slack}

On a polyhedron the minimizer lies on an unknown face. This paper gives no
deterministic polynomial-time method with a $\PosSLP$ oracle that finds this
face for general strongly convex quartics; Section~\ref{sec:constraints}
discusses that open case and several structured classes. The
difficulty disappears when the input bounds the nonzero optimal slacks away
from zero. Then ordinary approximation finds the full active set, and the
remaining problem is an unconstrained comparison on an affine subspace.

\begin{proposition}[Supplied slack gap]
\label{prop:upper-slack-gap}
Let the input consist of explicit sparse $f,h\in\Q[X_1,\ldots,X_n]$, a
rational $\mu>0$ with the promise $\nabla^2f\succeq\mu I$ on $\R^n$, a
rational polyhedron $P=\{x:Ax\le b\}$ with rows $a_i^{\mathsf T}$, a rational
$\delta>0$, and a relation ${\bowtie}$. When $P\ne\emptyset$, let $p_P$ be
the minimizer of $f$ on $P$, and assume the promise that every slack
$b_i-a_i^{\mathsf T}p_P$ is either $0$ or at least $\delta$. Then deciding
whether $P\ne\emptyset$ and $h(p_P)\bowtie0$ reduces to one $\PosSLP$
instance. The reduction also computes the set of active rows at $p_P$ by an
ordinary polynomial-time algorithm.
\end{proposition}

The proof (Appendix~\ref{app:upper-slack}) decides emptiness by linear
programming, obtains a feasible rational point whose distance to $p_P$ is at
most $\delta/(8\max_i\norm{a_i}_1)$ from convex value optimization, and
reads off the active set $S$ as the rows with slack below $\delta/2$. If
$d\in\ker A_S$, then $p_P\pm\theta d\in P$ for small $\theta>0$, so
$\nabla f(p_P)^{\mathsf T}d=0$; hence $p_P$ minimizes $f$ on the affine space
$\{x:A_Sx=b_S\}$. In a chart of this space with an identity block, the
restricted function keeps the modulus $\mu$, and the Newton machinery
applies. The chart is never expanded into monomials, so unary degree is
allowed. No multiplier bound, strict complementarity or Slater point is
needed. The slack gap is a promise that the algorithm does not check, and
it is a real restriction: Lemma~\ref{lem:separation} only guarantees nonzero
slacks of size $2^{-2^{\poly(L)}}$, which cannot be printed in polynomially
many bits. The proposition gives no deterministic $\mathrm P^{\PosSLP}$
bound for general polyhedra.

\subsection{Nested cube roots}
\label{sec:upper-roots}

The lower bounds of Section~\ref{sec:reductions} pass through a language of
certified cube-root circuits, in which every gate is the real cube root of an
affine combination of earlier gates and their squares, and supplied positive
intervals certify every gate. Its comparison problem has an upper bound by
the same scheme with two elementary substitutions. For a circuit with $k$
gates, multiplication by a common denominator turns the gate values into
algebraic integers in a field of degree at most $3^k$, and all conjugates of
the gate values are bounded by $2^{3L}$; the norm of the scaled difference
between the output gate $\xi_o$ and the threshold $r$ then separates
$\xi_o-r$ from zero by $2^{-2^{6L}}$. Newton's iteration for the cube root,
started at the upper end of the certified range, replaces the warm start. Theorem~\ref{thm:reductions-root-language} states the resulting
classification: strict and weak comparisons are $\PosSLP$-complete, and
equality has the upper bound only. The upper bound is a tailored instance of
the classical approximation-and-separation argument, not a new result about
general radical circuits.

\subsection{One sign test for adaptive sign tests}
\label{sec:upper-closure}

Several algorithms in this paper make many exact comparisons, and later
comparisons depend on earlier answers; examples are the exact quadratic
programming steps of Section~\ref{sec:constraints}. Such an algorithm is a
polynomial-time Turing reduction to $\PosSLP$. The next theorem converts it
into a reduction to one instance.

\begin{theorem}[Compiling sign tests]
\label{thm:posslp-closure}
\leavevmode
\begin{enumerate}[label=(\alph*)]
\item Let $C$ be a circuit with Boolean input nodes $u_1,\ldots,u_m$,
constants $0$ and $1$, binary $+$, $-$ and $\times$ gates, and unary threshold
gates $H(v)=1$ if $v>0$ and $H(v)=0$ otherwise. Let $S\ge1$ be its number of
nodes. Suppose that for every $u\in\{0,1\}^m$ all node values are integers
and the output is $0$ or $1$. In time polynomial in $S$, without oracle calls,
one can construct an integer circuit $C'$ of size $O(S^2)$ with the same
input nodes and no threshold gates such that $C'(u)>0$ exactly when
$C(u)=1$, for every $u\in\{0,1\}^m$.
\item Let $\Pi$ be a decision or promise problem decided by a deterministic
oracle Turing machine with oracle $\PosSLP$ that halts within a polynomial
number of steps on every input and for all oracle answers. Then $\Pi$
reduces to one $\PosSLP$ instance. Consequently, $\mathrm P^{\PosSLP}$ is the
class of problems that are polynomial-time many-one reducible to $\PosSLP$.
\item In particular, from integer circuits $C_1,\ldots,C_k$ and a Boolean
circuit $\beta$ with $k$ inputs one can construct in polynomial time one
integer circuit whose output is positive exactly when
$\beta([C_1>0],\ldots,[C_k>0])=1$.
\end{enumerate}
\end{theorem}

The proof is in Appendix~\ref{app:upper-closure}. In~(a), all exact values
have absolute value at most $2^{2^S}$, and distinct integers are at least
one apart. Each threshold gate is replaced by a rational map that is within
$2^{-2^{2S+4}}$ of the correct bit whenever its input is within $1/4$ of the
exact integer. The map first shrinks the magnitude from $[1,2^{2^{S+2}}]$ to
$[1,2]$ by the scaled steps $x\mapsto2\sigma x/(\sigma+x^2)$ with
$\sigma=2^{2^{S+2}},2^{2^{S+1}},\ldots,2^{2^1}$, and then applies
$x\mapsto2x/(1+x^2)$, the classical inverse Newton iteration for the sign
function, $2S+5$ times. An absolute-error induction over the original gates
shows that the errors never exceed $1/4$; all denominators are positive, so
the final pair $P/Q$ gives the integer output $2P-Q$. For~(b), the oracle
machine is unrolled into such a circuit. Each possible query is evaluated by
a padded universal interpreter: the query string is parsed by a Boolean
circuit, each instruction slot selects its operands by sums of earlier slot
values multiplied by address indicators, and malformed queries evaluate to
the answer no. The interpreter enumerates addresses, not answer transcripts,
so its size is polynomial although later queries depend on earlier answers.
Part~(c) is the special case of a nonadaptive query list.

The theorem changes the form of reductions; it gives no polynomial-time
algorithm for $\PosSLP$. It applies to deterministic computations with a
worst-case time bound. It does not remove randomness from Las Vegas
algorithms, does not turn the nondeterministic classes
$\mathrm{UP}^{\PosSLP}$ or $\mathrm{NP}^{\PosSLP}$ into deterministic ones,
and compiles one output bit at a time; a vector output such as an active set
is a polynomial list of separate instances. The exact separation of distinct
integers is essential, and the proof uses the full lengths of the compressor
and refinement schedules; no claim is made for shorter schedules.

The main application in this paper concerns constrained problems. For
strongly convex objectives on boxes with structured Hessians and for
separable objectives on network-flow polyhedra, Section~\ref{sec:constraints}
decides exact predicates at the constrained minimizer by deterministic
algorithms that solve exact quadratic subproblems on shared circuits and ask
adaptive sign queries. Their warm starts are ordinary computations, and the
number of later operations is bounded by a polynomial in the input length
that does not depend on the expanded sizes of the circuit values. Part~(b)
therefore compiles each such predicate into one $\PosSLP$ instance
(Corollary~\ref{cor:structured-single-sign}). The randomized
constraint-rank algorithm of Theorem~\ref{thm:constraints-rank} is not
covered, because its running time is bounded only in expectation.

\subsection{Relation to prior work}
\label{sec:upper-prior}

The combination of Newton approximation by short arithmetic circuits, an
algebraic separation bound and numerator-denominator division elimination is
established. Allender, B\"urgisser, Kjeldgaard-Pedersen and Miltersen showed
that Square Root Sum lies in $\mathrm P^{\PosSLP}$ by a polynomial-time
Turing reduction and characterized
$\mathrm P^{\PosSLP}$ as the Boolean part of constant-free polynomial-time
real computation \citep{AllenderEtAl2009}. Etessami, Stewart and Yannakakis
gave polynomial-time many-one reductions to $\PosSLP$ for threshold
comparisons of least fixed points of probabilistic polynomial systems with
the same architecture \citep[Appendix~C]{EtessamiStewartYannakakis2012}.
Exact semidefinite feasibility has comparable arithmetic-circuit reductions
\citep{TarasovVyalyi2008}. The warm start for strongly convex objectives uses
the elementary estimate $\norm p\le\norm{\nabla f(0)}/\mu$ to fix an explicit
box, and then applies the polynomial-time convex value interface of
Lemma~\ref{lem:convex-value}, which rests on the ellipsoid method.
For merely convex polynomial programs on unbounded domains,
\citet{SlotSteurerWiedmer2025} give radius and objective-gap approximation
bounds; see Section~\ref{sec:points}. The warm start for monotone maps uses
the rounded central-cut ellipsoid method of
\citet[Section~3.2]{GroetschelLovaszSchrijver1988}.

This section proves a one-instance upper bound for explicit globally strongly
convex objectives and strongly monotone maps of polynomially bounded degree,
for arbitrary polynomial observables and all six relations, including
equality, together with checked certificate languages. The lower bounds of
Section~\ref{sec:reductions} give matching $\PosSLP$ completeness for strict
and weak value and coordinate comparisons on certified quartics, and for
coordinate comparisons on their certified cubic gradient maps. Equality has
only the upper bound. The probabilistic polynomial systems of
Etessami, Stewart and Yannakakis form a different input class, and no
reduction between the two classes is used here. For Theorem~\ref{thm:posslp-closure}, the rational sign
iteration and its scaled variants are classical numerical tools, and the
oracle characterization of $\mathrm P^{\PosSLP}$ is known; the statement
proved here is the uniform compilation of adaptive integer sign tests into
one instance, with the error budget and interpreter given in full. This
compilation also converts the Square Root Sum oracle upper bound into a
one-instance upper bound. We make no priority claim for it.
